\section{Corrections found in this work}\label{app:corr}

This appendix lists the corrections made to the working draft (its sections on the killed kernel and on the Shell route), to formal statements and
to the preprint~\cite{DSouza}. The statements of this paper are the corrected ones; ``counterexample'' means an explicit one, and ``kernel-checked''
that it is a Lean theorem.

\subsection{Statements of the working draft that are false as written}
\begin{enumerate}
\item \emph{Lemma~\ref{lem:shell-1prime}, far part.} The draft allowed any cutoff $\Xi$ with $0\le\Xi\le1$ supported in $(-1,1)$, including $\Xi\equiv0$, for
which the far part is all of $a$ (numerical counterexample). Corrected by~\eqref{eq:Xi}; the constant becomes $2s/(\pi^2c\kappa)$. The same
hypothesis is needed in the model mass of Lemma~\ref{lem:shell-4} (written out in its proof).
\item \emph{Lemma~\ref{lem:shell-2}(c), the spike.} False without $\delta<1/(rQ)$; counterexample $a/r=0/1$, $\delta=1/Q$ for every $Q\ge2$. Not used.
\item \emph{Theorem~\ref{thm:shell-K} as stated in the draft.} The kernel on $(\log Q+4,\Lc]$ is $(\Lc+\log K+O(1))/s$, not $\abs\alpha$; the pointwise claim beyond
$\alphap$ fails on the top edge $(\log X-1,\log X]$; ``by the symmetry $s\mapsto-s$'' is not a symmetry of $\mathbf F$. The corrected statement is integrated and uses
the weight $\omega_K$. None of this changes the constant or the rate.
\item \emph{The condition on $\epsilon$.} The draft required only $\epsilon\to0$; Lemma~\ref{lem:shell-6d} needs $\log(1/\epsilon)=O(\log\log Q)$, through $R_1$ in
$\mathcal Z$. With $\epsilon=Q^{-c'}$ the bulk condition at $\alphap=2497/1500$ holds only for $c'<0.0012$. The draft's choice $\epsilon=1/\Lc$ is fine.
\item \emph{The rate.} The draft stated the error $(\log\log Q)^4(\log Q)^{-503/1994}$; $503/1994$ is the value at $\eps''=0$, and the proof gives
$(\log Q)^{-\theta}$ for every $\theta<503/1994$.
\item \emph{Lemma~\ref{lem:shell-6c}.} ``Every block has $R\ge1$'' holds only for nonempty blocks, and the cover needs $R_1\le N_0$ (numerical control).
\item \emph{Lemma~\ref{lem:shell-5c}(2).} Needs $T\le N_0$; without it no absolute constant exists (numerical). Part (3) with the Weil-form error needs the
actual size of $R_1$, not only $R_1\le Q/(2K)$.
\end{enumerate}
Further items concern formal transcriptions, not the draft. The formal statement of Lemma~\ref{lem:shell-6b} over the reals needs $j\in\mathbb N$
and $H\ge0$, which were added as hypotheses (counterexamples $j=-3$; $H=-100$ with real powers). Three further false statements: the first Lean form
of Lemma~\ref{lem:K-G} put the last term inside the integral (kernel-checked refutation); the first Lean form of Lemma~\ref{lem:shell-2}(a) was
stated for all $Q\ge2$ and is false at $Q=2$, where $H_{\mathfrak w}=0$, so the lemma is stated in eventual form; and the first Lean form of
Lemma~\ref{lem:shell-star} (\texttt{ring\_le\_primitive}) was false as Lean parses it. In that form the body of the last integral
extends to the end of the expression and takes in the trailing term $+2(\log K+C_\mu)U$, which so sits inside the integral and inside both sums (the
parse is a Lean theorem, \texttt{frozen\_rhs\_parse}, proved by \texttt{rfl}); with $R_1=1$, $K=3$, $C_\mu=0$, $Q=1000$, $a_n=e(-2n/Q)$ on the primes
$1000<n\le30000$ and $\mathcal M=S$, all hypotheses hold and the ring mass is $331.4$, against $59.4$ for the right side as parsed (numerical
counterexample). The corrected statement \texttt{ring\_le\_primitive\_corr} brackets the integral, and the first form was removed from the library.
This is a correction of a Lean statement, not of the paper: the inequality displayed in Lemma~\ref{lem:shell-star}, which has the term outside the
integral and the sums, was never affected. Of the items above, the far part of Lemma~\ref{lem:shell-1prime}, the transcription of
Lemma~\ref{lem:K-G}, the first form of Lemma~\ref{lem:shell-2}(a) and the first form of Lemma~\ref{lem:shell-star} were found by the formalisation,
the others in the proofs on paper.

\subsection{Missing hypotheses, constants and wording}
\begin{enumerate}
\item Lemma~\ref{lem:shell-2}(a) needs $R_1\ge1$; its closing sentence needs $N/\epsilon\le Q^2(\log Q)^{-12}$; the per-class error in Step~2 is $2V^*$, not $V^*$; the
tail bound of Step~3 could not be reproduced and is replaced by one with the same conclusion; Step~4 holds up to $O(\delta(1+\log Q))$.
\item Lemma~\ref{lem:shell-P}(iv): the bound $\sum_e\abs{c_e}2^{\omega(e)}/e\le\zeta(2)^2/\zeta(4)$ holds for the plain kind and fails for the $q/\varphi$ kind ($3.2796$);
only convergence is used.
\item Lemma~\ref{lem:shell-6a} needs $K\le Q$ and $T\le N_0$ for its tail step, and $N_0\ge e^3$, which the ``large $N_0$'' of its proof and its
formal statement use.
\item Lemma~\ref{lem:shell-1prime}: the prime-power peak is $\ll T^2(\log N)^2N^{-1/2}$, not $T(\log N)^2N^{-1/2}$.
\item Lemma~\ref{lem:shell-star}: the constant $1.4708$ is numerical only; any $O(1+\log K)$ bound suffices.
\item Proposition~\ref{prop:shell-Z}: the explicit formula is used in the Weil form, with the error $E'$ in place of $E$. The draft's
remark that the explicit formula for Dirichlet $L$-functions is not in the Lean artifact of~\cite{AF} is out of date. Lemma~\ref{lem:shell-5c}(1): ``total
$\ll T$'' does not name the norm.
\item Theorem~\ref{thm:shell-S} is stated with $\norm W_{C^A}Ts_0$ on the right, with $K$ and $\epsilon$ explicit; the draft's
$\int W_j\norm a^2\asymp\gwin(j)Ts_0$ in the transfer would need a lower bound on the mass of $W_j$.
\item The working draft's remark that Tao's inequality can be avoided needs $\epsilon\log Q\to0$, which does not hold for $\epsilon=1/\Lc$; the ring points
have $q>Q/K$, which is enough.
\item Lemma~\ref{lem:K-S} is proved only for $s\le\log X$, the range used; the draft stated it for $N\ge Q$ (D3).
\item Lemma~\ref{lem:K-P}: the draft's prime number theorem with error $y\exp(-c\sqrt{\log y})$ is replaced by the weaker~\eqref{eq:PNT}, and its unverified
Mertens constant $\abs{R(x)}\le2$ by $\log4+4$, so the bounds $E_2^2/U\le33.3\,\lambda\Lc/T$ and $2E_2/\sqrt U\le16.4(\Lc/T)^{1/2}$ become
$E_2^2/U\le86\,\lambda\Lc/T$ and $2E_2/\sqrt U\le26.3(\Lc/T)^{1/2}$ (D4). The step $\sum_n\abs{\tilde{\mathfrak m}(n)}^2\ge U-1$ was asserted
without proof and is written out here.
\item Lemma~\ref{lem:K-Z}(b) holds for all $r,r'\le Q$; the draft's $r+r'\le Q$ is more than needed.
\item Theorem~\ref{thm:one-prime}: the draft's $0.72375738$ is at $\lambda=32/25$, outside the formal design, and its statement that a profile at $\lambda=5/4$ ``fits the
caps'' does not hold, because the formal design fixes $\lambda=\lambda^*$ (D1); the design profile of~\cite{DSouza} already gives $0.7235$ (D2); a table entry was rounded
rather than truncated (D5).
\item The bandwidth: the draft's open item on the uniformity in $\lambda$ of the rows of~\cite[\S\S7--10]{DSouza} is resolved in \S\ref{ssec:bandwidth} for $\FQ$ at
$\lambda=191/100$ (the dyadic family at $\lambda=93/50$ was not examined, and its constant is not claimed; Remark~\ref{rem:shell-dyadic}); its list of
formal caps was incomplete; the ends constant becomes $c_1=0.52$ with the Minkowski split; the moment floor stays at $3/4$ in eventual form; the ramp rate is
$\Theta(1/\log Q)$; the zone slope is $1.9458$ with certified Lipschitz constant $1.906$ (a doc-string said $1$); the Gevrey gate is raised to $5\cdot10^4$.
\item Lemma~\ref{lem:shell-6d}: the near bound holds, and is proved, for $0\le X_0\le s_0/5$, where the near zeros have $\beta\ge\frac45$ and (LF) counts them.
The draft stated it with no range of $X_0$, and the first Lean form \texttt{Z6d\_counts} for every $X_0\ge0$; for $X_0>s_0/5$ the near sum contains zeros
with $\beta<\frac45$, and (M) gives only a larger exponent there (it would need $\eps''\ge1-1/\alphap\approx0.399$). The corrected formal statement
\texttt{Z6d\_counts\_corr} adds $5X_0\le s_0$ to the near clause only and is proved; the first form is kept, open and unused.
Theorem~\ref{thm:shell-S} uses the bound at $X_0=O(\log\log Q)$ only, which its proof checks. No constant or rate changes.
\item The pointwise assembly~\eqref{eq:shell-assembly}, formal statement. The first Lean form (\texttt{AS\_pointwise}) allowed any $\alpha''<2$,
independently of $\lambda$. Its ring is that of the near polynomial not cut at $X$, while $\mathbf F(s)$ sees $a(s)$ cut at $X$; for $s+1>\lambda\Lc=\log X$
the argument of the assembly bounds the ring of the cut polynomial, which that ring does not control. No counterexample is known. The corrected form
\texttt{AS\_pointwise\_corr} adds $\alpha''<\lambda$ and is proved; the first form is kept, open and unused. Theorem~\ref{thm:shell-K}
uses the formal assembly with $\alpha''<5/3<191/100=\lambda$, and its formal statement is unchanged. In this paper the ring is that of $a(s)$, which is cut
at $X$ (\S\ref{ssec:notation}), and the assembly is used only for $s\le\alphap\Lc$ with $\alphap<\lambda$, where the cut and the uncut near polynomials
agree, as they do in Theorem~\ref{thm:shell-S} through the condition $s_0\le\log X-2$. In \S\ref{ssec:shell-notation} and
in the assembly the range of $s$ also carries $s\le\log X-1$, which $\norm{a(s)}^2=U(s+O(1))$, Lemma~\ref{lem:shell-4} and the
relative costs of Lemma~\ref{lem:shell-1prime} need, and which $\alpha''<\lambda$ implies for large $Q$. No constant or rate changes.
\end{enumerate}

\subsection{\texorpdfstring{Items for the preprint~\cite{DSouza}}{Items for the preprint}}
\begin{enumerate}
\item Its formalisation section says that the six headline theorems depend on one \texttt{sorry}, the sharp large-sieve constant. The large sieve
is now proved in Gallagher's form $Q^2+\pi N$, and the six theorems are kernel-checked with no \texttt{sorry}; the budget constant
of~\cite[Lemma~4.3]{DSouza} becomes $1+4Q^{-\delta}$ instead of $1+2Q^{-\delta}$.
\item The sieve-budget row $L_{11}$ should read $Q^{\lambda-2}(T/2\pi)^\lambda$.
\item The even and odd subfamilies now have the constants $0.7212$ and $0.7098$ of the full families, by the reflected large sieve; this is a Lean theorem on a
branch of the formal tree, and it supersedes the preprint's $0.6980$ and $0.6919$.
\item The bibliography of~\cite{DSouza} gives a different author for the source of H1--H6; this paper cites the arXiv listing~\cite{AF}.
For reference: the preprint's own text says of that source that it is not posted to a preprint server~\cite[\S1.2]{DSouza}; the arXiv paper~\cite{AF}
states on its first page that its mathematical argument was discovered and written by Claude, an AI developed by Anthropic, and that the listed
authors verified the proof and take responsibility for its content [V].
\end{enumerate}
